
where the vertical arrows are the canonical homomorphisms @@M0@@ (IV 2.5.1).

\sgaStableTarget{sga6.e10.sec.7.7}\subsection*{7.7. Specialization for @@M1@@ and @@M2@@.}
When @@M3@@ and @@M4@@ are regular (for example, when @@M5@@ is smooth), the two preceding vertical arrows are isomorphisms (IV 2.5), and the specialization homomorphism can then be interpreted as a homomorphism
\sgaStableTarget{sga6.e10.eq.7-7-1}@@M6@@
which can also be described as the unique homomorphism making the following diagram commutative
\sgaStableTarget{sga6.e10.eq.7-7-2}@@M7@@
Since, in this last diagram, @@M8@@ and @@M9@@ are ring homomorphisms (IV 2.7 b), with @@M10@@ surjective, it follows that \emph{the specialization homomorphism} (7.7.1) \emph{is a ring homomorphism}. For the same reason, \emph{if} @@M11@@ \emph{is separated}, \emph{the homomorphism} (7.7.1) \emph{is even a homomorphism of} @@M12@@-\emph{rings}. To make this precise, one need only note that a regular separated noetherian scheme is divisorial (II 2.2.7.1), so its @@M13@@ is identified with its ``naive'' @@M14@@ (IV 2.9 (ii)), which is naturally endowed with a @@M15@@-ring structure (VI 3.2). On the other hand, the regularity of @@M16@@ and @@M17@@ implies that of @@M18@@ (EGA @@M19@@ 17.1.8), and the homomorphisms @@M20@@ and @@M21@@ are @@M22@@-homomorphisms, so the same is true of (7.7.1), which is obtained from @@M23@@ by passage to the quotient.
